\section{Wavelet-Based Image Coding}




\begin{questionbox}{Domanda 29}
State the time-frequency uncertainty principle ($\Delta t \cdot \Delta f \geq 1/4\pi$), explain why it imposes a trade-off, and how wavelets address it with adaptive multiresolution (short windows at high frequencies, long at low frequencies).
\end{questionbox}



$$
\Delta t \cdot \Delta f \geq \frac{1}{4\pi}
$$


\begin{itemize}
  \item Short window → good time/space localization, poor frequency resolution.
  \item Long window → good frequency resolution, poor localization.
  \item Wavelets → multiresolution:
  \item high frequencies → short basis functions;
  \item low frequencies → long basis functions.
  \item Image match → edges need localization; smooth trends need coarse-scale representation.
\end{itemize}



\begin{questionbox}{Domanda 30}
Compare STFT (rigid tiling) and \textbf{DWT} (adaptive tiling) of the time-frequency plane, linking them to the trends vs anomalies image model.
\end{questionbox}


\begin{itemize}
  \item \textbf{STFT} → fixed window size → rigid time-frequency tiling; same resolution for all frequencies.
  \item \textbf{\textbf{DWT}} → filter banks + downsampling → low-frequency approximations + high-frequency detail subbands.
  \item Repeated \textbf{DWT} → multiresolution representation.
  \item Image model:
  \item smooth trends → coarse low-frequency components;
  \item anomalies/edges → high-frequency localized details.
\end{itemize}



\begin{questionbox}{Domanda 31}
(TC\&JPEG) What is a primary advantage of the hierarchical (multiresolution) decomposition offered by the \textbf{wavelet} Transform in JPEG2000?
\end{questionbox}


\begin{itemize}
  \item \textbf{wavelet} decomposition → progressive and multiresolution coding.
  \item Low-resolution approximation decoded first → detail subbands refine quality and resolution.
\end{itemize}



\begin{questionbox}{Domanda 32}
(TC\&JPEG) Compare the block-based \textbf{DCT} approach used in JPEG with the \textbf{wavelet}-based decomposition used in JPEG2000, specifically in terms of how they handle the image signal and the resulting artifacts.

(TC\&JPEG) Which of the following is a key reason why JPEG2000 is generally more efficient than JPEG?
\end{questionbox}


\begin{itemize}
  \item \textbf{JPEG}:
  \item independent $8\times8$ block \textbf{DCT};
  \item simple and efficient;
  \item low bitrate → \textbf{blocking artifacts}.
  \item \textbf{JPEG2000}:
  \item \textbf{DWT} over large tiles / whole image;
  \item multiresolution representation;
  \item embedded bitplane coding and precise \textbf{rate} control;
  \item low bitrate → ringing/blurring near edges, not blocking.
  \item Efficiency reasons:
  \item \textbf{wavelet} multiresolution decomposition;
  \item embedded bitstream truncation;
  \item independent codeblocks;
  \item quality/resolution scalability;
  \item no fixed $8\times8$ boundaries.
\end{itemize}



\begin{questionbox}{Domanda 33}
Describe the JPEG2000 architecture (Tier 1: \textbf{DWT} 9/7 or 5/3 -> fine \textbf{quantization} -> arithmetic coding of codeblocks per bitplane; Tier 2: \textbf{EBCOT}). Where does the \textbf{lossy} operation actually happen?
\end{questionbox}


\begin{itemize}
  \item \textbf{JPEG2000} → \textbf{wavelet} still-image codec for scalability, \textbf{rate} control, ROI, \textbf{lossy}/\textbf{lossless}.
  \item Tier 1:
\end{itemize}


\begin{figure}[H]
\centering
\begin{tikzpicture}[
    node distance=1.1cm,
    block/.style={draw, fill=blue!5, text width=5cm, align=center, minimum height=0.6cm, rounded corners=3pt, font=\small},
    arrow/.style={thick,->,>=stealth}
]
    \node (a) [block] {Image};
    \node (b) [block, below=of a] {DWT};
    \node (c) [block, below=of b] {Quantization};
    \node (d) [block, below=of c] {Codeblocks};
    \node (e) [block, below=of d] {Bitplane Arithmetic Coding};

    \draw [arrow] (a) -- (b);
    \draw [arrow] (b) -- (c);
    \draw [arrow] (c) -- (d);
    \draw [arrow] (d) -- (e);
\end{tikzpicture}
\caption{JPEG2000 Tier 1 Coding Chain}
\end{figure}


\begin{itemize}
  \item 9/7 \textbf{DWT} → irreversible, \textbf{lossy}.
  \item 5/3 \textbf{DWT} → reversible, \textbf{lossless}.
  \item Coefficients → codeblocks → coded bitplane by bitplane.
  \item Tier 2 / \textbf{EBCOT} → layers, packets, progression orders, truncation points.
  \item \textbf{lossy} operations → \textbf{quantization} and embedded bitplane/pass truncation.
\end{itemize}



\begin{questionbox}{Domanda 34}
Compare JPEG and JPEG2000 on channel-error robustness (codeblock independence, contained vs catastrophic propagation, resynchronization markers).
\end{questionbox}


\begin{itemize}
  \item \textbf{JPEG}:
  \item sequential \textbf{entropy} coding;
  \item bit error can desynchronize Huffman stream;
  \item corruption propagates until restart/resynchronization marker.
  \item \textbf{JPEG2000}:
  \item relatively independent codeblocks;
  \item stream organized into packets/layers;
  \item damage tends to remain local to codeblock/\textbf{subband}/quality layer.
  \item Result → JPEG2000 more robust to localized channel errors.
\end{itemize}

---


